% !TeX root = ../main.tex
\section{Data Collection}
\label{sec:data}

Our training corpus combines real-world data collection, captures from a simulated environment, and open-source datasets (Fig.~\ref{fig:data_pipeline}).
The corpus contains 600,000 clips, including 200,000 simulation clips, and provides complementary observations, scene diversity, and geometric supervision.

\subsection{Data Sources}
\label{sec:data_sources}

\paragraph{Real-world data collection.}
We record onboard videos during real-world UAV flights, capturing observations under physical flight conditions.
We supplement these recordings with YouTube drone videos to expand the diversity of real-world aerial footage.

\paragraph{Simulation data collection.}
We build an automated collection pipeline in GTA V with DLSS 5 neural rendering~\cite{nvidia2026dlss5} to capture first-person UAV flights close to the ground.
An obstacle-avoiding flight controller supports continuous collection, while randomized turns diversify the trajectories.
We jointly export video, real-time camera poses, and linear depth maps derived from the rendering engine's depth buffers and aligned with the enhanced RGB frames.
The annotations share the simulator's approximately metre-scale units.
The combination of realistic imagery and accurate camera and depth annotations makes this our highest-quality source for geometric supervision.

\paragraph{Open-source datasets.}
We use Sekai-Walking-Clean and SpatialVID-Clean from SolarWM~\cite{huang2026solarwmopendatascalable}, derived from Sekai~\cite{li2025sekaivideodatasetworld} and SpatialVID~\cite{wang2025spatialvidlargescalevideodataset}, retaining only the \texttt{high} and \texttt{xhigh} quality tiers.
These clips are filtered and processed using distilled LTX-2.3-22B with the Clean Plate IC-LoRA adapter to remove people and vehicles.
We also include the drone subset of Sekai for additional aerial observations.
% Model identity verified against SolarWM Sec. 5.3 and its released manifest:
% https://arxiv.org/html/2609.02886v1#S5.SS3
% https://github.com/Enderfga/solarwm-data/blob/main/configs/clean_plate_models.json

From DL3DV-10K~\cite{ling2024dl3dv10klargescalescenedataset}, we reconstruct scenes with 3D Gaussian Splatting (3DGS) and render videos along designed camera trajectories.
We apply Difix3D+~\cite{wu2025difix3d} to repair dark patches and streaking artifacts, then use optical-flow correspondences between successive frames for temporal smoothing to reduce flicker.

\subsection{Camera Annotation}
\label{sec:camera_annotation}

For videos without camera labels, Pi3X~\cite{wang2026pi3permutationequivariantvisualgeometry} initializes poses and scene geometry, followed by bundle adjustment (BA) with shared intrinsics.
We align UniDepthV2~\cite{piccinelli2025unidepthv2} predictions with reliable reconstructed points and apply bounded inverse-depth correction to recover one positive scale per clip, preserving rotations and trajectory shape.
Section~\ref{sec:annotation_evaluation} presents the selection evidence, and the Supplementary Material gives implementation details.

\subsection{Depth Annotation}
\label{sec:depth_annotation}

For non-GTA videos, MetricAnything Student PointMap~\cite{ma2026metricanything} provides approximately metric optical-axis depth.
Metric Video Depth Anything (VDA) Large~\cite{chen2025videodepthanything} stabilizes framewise scale while preserving spatial structure.
We temporally stabilize sky masks from Depth Anything 3~\cite{lin2025depthanything3}, then downsample depth and apply the shared logarithmic encoding in Sec.~\ref{sec:bidirectional_training}.
Section~\ref{sec:annotation_evaluation} summarizes GTA validation, and the Supplementary Material details alignment and sky processing.

% Depth implementation verified against AutoDL's deployed metric_depth_models.py,
% metric_depth_processing.py, metric_depth_contract.py, and source-config.local.json
% under df-full-20261102, plus the GTAAligned depth-cache metadata.
